\section{The actual little Julia arcs and KKL ideal boundaries}
\label{sec:marking}
All markings in this section refer to the \emph{physical}
slit complements of the proper central restriction.
They are not a substitute for a marked ambient puzzle
or a specification of external ray angles.

\begin{theorem}[Physical period-three marking, conditional on (H1)--(H2)]
\label{thm:marking}
Assume the geometric renormalization and physical ideal-marking
bridges \emph{(H1)}--\emph{(H2)} of
Definition~\ref{def:bridges}.
Put $t=x_{6p}=-x_{3p}>0$, and
\begin{equation}\label{eq:tree}
 K_0=[-t,t],\quad K_1=[x_{4p},x_p],\quad
 K_2=[x_{2p},x_{5p}],\quad
 \T_p=[x_{2p},x_p],\quad
 \T'_p=[-x_p,x_p].
\end{equation}
The renormalization of $G_p$ of period three is
\emph{prime and primitive}, and its \emph{actual}
little filled Julia sets are $K_0,K_1,K_2$.
The KKL Hubbard tree is exactly $\T_p$, with
dividing fixed point $\alpha_p$.
Moreover $\T'_p=G_p^{-1}(\T_p)\cap V_p$.

The target marked surface $\Om_p=U_p\setminus\T_p$
has ten consecutive \emph{inner} ideal atoms
$I_0,\ldots,I_9$, indexed by the five physical
increasing cuts
\[
 (a,b,d,e,h)=(x_{5p},x_{3p},x_{6p},\alpha_p,x_{4p}).
\]
The source $\Om'_p=V_p\setminus\T'_p$
has twenty consecutive inner ideal atoms
$S_0,\ldots,S_{19}$ indexed by
\[
 (-x_{4p},x_{2p},x_{5p},-\alpha_p,x_{3p},
 x_{6p},\alpha_p,-x_{5p},-x_{2p},x_{4p}).
\]
The two \emph{outer} ideal circles each have one
whole-component atom, not ten or twenty.
The unique new source-tree attachment is
$\T'_p\setminus\T_p=[-x_p,x_{2p})$, with
consecutive source sides $S_{19},S_0,S_1$.
\end{theorem}

\begin{proof}
The exact equation $x_{6p}=-x_{3p}$,
evenness of $G_p$ and of its iterates,
and \eqref{eq:order1} give
\begin{equation}\label{eq:portrait}
 G_p^3(0)=-t,\quad G_p^3(-t)=G_p^3(t)=t,\qquad
 G_p(K_0)=K_1,\quad G_p(K_1)=K_2,\quad
 G_p(K_2)=K_0.
\end{equation}
The full third-return disk in Theorem~\ref{thm:rouche}
has a central proper degree-two restriction.
In the normalized coordinate $w=A_pz$ its
connected filled Julia set contains the actual
real invariant interval $[-|A_p|t,|A_p|t]$:
the interval is forward invariant by \eqref{eq:portrait},
lies in the certified central disk component and
contains the sole critical point.  The Douady--Hubbard
connectedness criterion
\cite[Proposition 2]{DH1985} therefore makes the
filled Julia set connected.  By the
\emph{Straightening Theorem}
\cite[Theorem 1]{DH1985} it is hybrid-conjugate
to a unique monic centered quadratic $z^2+\gamma$.
Its critical portrait, in the normalized orientation,
is $0\mapsto u\mapsto -u\mapsto -u$ with $u>0$.
For a monic quadratic the relation
$f_\gamma^3(0)=f_\gamma^2(0)$ reduces
\emph{exactly} to $\gamma^3(\gamma+2)=0$.
The nonfixed critical point excludes $\gamma=0$,
so $\gamma=-2$.  Since $K(z^2-2)=[-2,2]$,
the filled Julia set of the third return is an arc.
It contains the entire straight segment joining
its endpoints $\pm u$, and an arc cannot contain
a proper different arc with the \emph{same}
endpoints.  Thus it is exactly this segment,
not merely homeomorphic to it.  Rescaling and
\eqref{eq:portrait} yield $K_0,K_1,K_2$ as
the three actual little filled Julia intervals.

The strict separations in \eqref{eq:order1}
make the three little filled sets pairwise
disjoint; this is \emph{primitive}
in \cite[\S2.3]{KKL2026}.  We prove primeness
rather than infer it from the word primitive.
The global polynomial $f_{c_p}$ is postcritically
finite and its critical point is strictly
preperiodic: $f^{9p}(0)=f^{6p}(0)=t\ne0$,
whereas $f^{3p}(0)=-t\ne0$.
The Misiurewicz local-connectivity/dendrite and
periodic-repulsion theorem
\cite[Theorem 27.4]{LyubichNotes}
and the subhyperbolic treatment
\cite[\S17]{Milnor2006} make its connected
Julia set a locally connected dendrite,
with all periodic cycles repelling.
The \emph{global} filled Julia sets agree,
$K(G_p)=K(f_{c_p})$, since $G_p=f_{c_p}^p$;
they are connected because the \emph{unique
critical point of the quadratic polynomial
$f_{c_p}$} has bounded orbit.  This does
not use the false premise that the
degree-$2^p$ global polynomial $G_p$
has only one critical point.  The relevant
central little filled set, once its
connectedness is proved for the restricted
proper map, is a subcontinuum of the
ambient dendrite and hence is a dendrite.

Suppose a connected period-two degree-two
renormalization of $G_p$ had little filled set
$L\ni0$.  Since the even-return critical orbit
lies in $L$, both $x_{2p}$ and $x_{4p}$
belong to it.  Uniqueness of dendrite arcs
then forces $[x_{2p},x_{4p}]\subset L$.
By the real monotonicity and \eqref{eq:portrait},
$L\cap G_p(L)$ contains the nondegenerate
interval $[x_{5p},x_{4p}]$.
Under the \emph{nestedness/straightening
transport} clause of (H1), these are
little filled Julia sets to which
McMullen's \emph{almost-disjoint little Julia
sets theorem} \cite[Theorem 7.3]{McMullen1994}
applies, so distinct such sets meet in
at most one repelling periodic point;
hence $L=G_p(L)$.  The theorem is
\emph{not} applied directly to the
unrestricted degree-$2^p$ polynomial.
But $0\in G_p(L)$ then has an inverse
$y\in L$, $y\ne0$ because $G_p(0)=x_p>0$.
The proper-domain containment demanded
by (H1) places both $0$ and $y$ in
the alleged degree-two return domain;
there they are distinct critical points
of $G_p^2$, contradicting its degree.
Without that domain containment the
original argument would be incomplete.
Thus no period-two renormalization exists.
The period-three renormalization is prime.
(The period-one case is the original $G_p$.)

Real invariance and the order \eqref{eq:order1}
show $\T_p=[x_{2p},x_p]\subset K(G_p)$ and
$G_p(\T_p)\subset\T_p$.  The three little
Julia intervals have extreme endpoints
$x_{2p}$ and $x_p$.  By dendrite arc
uniqueness any connected subset containing
them contains their real joining interval.
KKL's \emph{smallest connected Hubbard
continuum} definition
\cite[\S2.5]{KKL2026} therefore gives
precisely $\T_p$, not a conjectural drawn
subtree.  The unique fixed point
$\alpha_p\in(x_{6p},x_{4p})$ has the
dividing property by real cut order, as
in KKL \S2.4.  Finally
$G_p([-x_p,x_p])=\T_p$
under the two strictly monotone sheets.
The proper degree-two count exhausts
the two central inverse sheets; therefore
$G_p^{-1}(\T_p)\cap V_p=[-x_p,x_p]$,
not an additional off-real preimage arc.

For the \emph{ideal boundaries}, $U_p\setminus\T_p$
is a doubly connected domain with a round
outer boundary and a locally connected real
slit inner boundary.  So is
$V_p\setminus\T_p'$: the central proper degree-two preimage
component has a finite analytic
lemniscate boundary and hence locally
connected boundary (possible critical
boundary values require care);
a Jordan boundary is \emph{not}
inferred merely from the Rouch\'e count.
Also $\T'_p\Subset V_p$.
Cut either annulus along a short crosscut
joining its two boundary components.
The remaining simply connected slit
domain has locally connected boundary.
The Carath\'eodory prime-end
compactification \cite{Caratheodory1913}
provides an ideal circle for each
doubly-connected slit complement;
the further \emph{physical identification}
of all accesses with the claimed KKL
labels is precisely condition (H2).
Gluing the two crosscut banks then recovers the \emph{two}
ideal boundary circles of KKL Appendix A.1.
An interior point of the real slit has
exactly two local approach sectors
(above and below the slit) in a small
disk, and a terminal slit tip has
exactly one.  Shrinking crosscuts
show no third access is hidden.
The five inner target cut positions
therefore give ten inner ideal atoms,
and the ten source cut positions give
twenty.  The \emph{outer} circle
is entirely uncut, hence contributes
one ideal atom.  The physical order
\eqref{eq:order2} identifies all source
cuts.  These are actual
Carath\'eodory/KKL slit-side atoms,
without an unproved external ray-angle
atlas or an extra family-dependent
marking hypothesis.  Since $\T'_p$
extends $\T_p$ only to the left,
the set difference is the asserted
\emph{single} attachment, with
sides $S_{19},S_0,S_1$ by cyclic order.
\end{proof}

\begin{theorem}[Covering, immersion and incidence, conditional on (H1)--(H2)]
\label{thm:incidence}
Assume \emph{(H1)}--\emph{(H2)} of
Definition~\ref{def:bridges}, including the KKL
prime-end and pulled-marking compatibility.
On $\Om'_p$ the proper unbranched map
$G_p:\Om'_p\to\Om_p$ induces the
\emph{cover} $G_{p*}$, while the
nonproper inclusion $\iota:\Om'_p\hookrightarrow
\Om_p$ induces the \emph{set-valued proper-path
restriction} $\iota^*$.  With the above cyclic
indices their exact distinct tables are
\begin{equation}\label{eq:cover}
 G_{p*}(S_s)=I_{s\bmod10},\quad 0\le s<20,
\end{equation}
\begin{equation}\label{eq:blocks}
\begin{array}{c|l@{\qquad}c|l}
I_j&\iota^*I_j&I_j&\iota^*I_j\\ \hline
I_0&S_{18},S_{19},S_0,S_1,S_2&
I_5&S_{10}\\
I_1&S_3,S_4&I_6&S_{11},S_{12},S_{13}\\
I_2&S_5&I_7&S_{14}\\
I_3&S_6&I_8&S_{15}\\
I_4&S_7,S_8,S_9&I_9&S_{16},S_{17}
\end{array}.
\end{equation}
The attachment makes the $I_0$ block
not fully proper.  For the specified
crossing-attachment relative isotopy
class the two proper source pairs are
$(S_3,S_1)$ and $(S_{19},S_{17})$;
for the avoiding class the one pair
is $(S_3,S_{17})$.  Their inclusive
short source spans are $3,3,7$;
their \emph{cover-pushed} target
pairs are $(I_3,I_1),(I_9,I_7),(I_3,I_7)$,
with short spans $3,3,5$.
All pass the KKL A.1 filter $>2$.
No arbitrary A.1 parent is asserted
to have two restrictions.

For $E_0=(x_{5p},x_{3p})$,
$E_1=(x_{6p},\alpha_p)$,
$E_2=(\alpha_p,x_{4p})$,
the image incidence, indexed by output
rows and input columns, is
\begin{equation}\label{eq:M}
 M=\begin{pmatrix}0&0&1\\1&0&1\\1&1&0\end{pmatrix},
 \quad\det(zI-M)=(z+1)(z^2-z-1),\quad
 \rho(M)=\varphi=\frac{1+\sqrt5}{2}.
\end{equation}
The positive left Perron vector is
$(\varphi,1,\varphi)$.
Target zero-old-edge atom types are
$I_0,I_2,I_5,I_8$.
The atomic parent crest $I_0$ and
the cover-pushed first-child crests
$I_2,I_8$ have no old-edge carrier.
The two banks $I_1,I_9$ are \emph{one}
physical edge $E_0$, not two edges.
\end{theorem}
\begin{proof}
The restriction $G_p$ has no
critical points in $\Om'_p$:
its only central critical point
$0$ was removed with $\T'_p$.
The degree-two properness of
Theorem~\ref{thm:rouche} makes
$G_p:\Om'_p\to\Om_p$ an unbranched
cover.  KKL Appendix A.3.1
extends it to the ideal circle,
where one rotation-compatible
sheet labeling gives
\eqref{eq:cover}; every target
atom has two preimages.

By contrast the \emph{inclusion}
$\Om'_p\hookrightarrow\Om_p$
is nonproper precisely because
the source slit extends to the
left of the target one.  KKL
\cite[Proposition 6.5]{KKL2026}
describes the pullback of an
ideal interval as inverse-side
accesses \emph{together with}
Carath\'eodory boundaries of
attached tree components.
Traversing the physical
ten source cuts in
\eqref{eq:order2} and inserting
the three attachment sides
produces successively the
ten consecutive source blocks
\eqref{eq:blocks}, which
partition all twenty indices
exactly once.  The attached
part of the $I_0$ block
maps by inclusion to target
\emph{interior}; no point map
$I_0\mapsto S_s$ can replace
$\iota^*$ here.  KKL
\cite[Proposition 6.6 and
Appendix A.3.2--A.3.4]{KKL2026}
then gives the proper-path
restriction semantics.

A crossing arc whose right
intersection belongs to
$(x_{5p},-\alpha_p)$ and
whose left intersection belongs
to $(-x_{4p},x_{2p})$ splits
at the attached tree into two
proper source components, whose
upper/lower endpoints, read from
the cut order, are $(3,1)$ and
$(19,17)$.  If the specified
comparison arc avoids the attached
piece and its left endpoint is
left of $-x_p$, it leaves the
one proper pair $(3,17)$.
For twenty cyclic atoms the
inclusive short lengths are,
respectively, $3$, $3$ and $7$;
for example the shorter arc
from $S_3$ to $S_{17}$ is
$S_{17},S_{18},S_{19},S_0,S_1,S_2,S_3$.
Applying the \emph{different}
cover table gives target spans
$3,3,5$.  Every one exceeds two,
so none is removed by KKL A.1.
This is a statement about these
\emph{relative isotopy classes},
not a selection theorem for all
A.1 Waves.

The five cut-end images under
$G_p$ have index sequence
$(2,4,4,3,0)$, by the exact
orbit identity and fixed
point.  On the old real edges
$(a,b),(d,e),(e,h)$,
coverage counts give:
$E_0$ maps across $E_1,E_2$,
$E_1$ across $E_2$,
and $E_2$ across $E_0,E_1$.
This is exactly \eqref{eq:M}.
Its characteristic factorization
is a direct determinant.
Since $M^4$ has strictly
positive entries, Perron--Frobenius
\cite[Chapter XIII, \S2]{Gantmacher1959}
identifies the positive eigenvalue
$\varphi$ as the simple spectral
radius; alternatively
$(\varphi,1,\varphi)M
=\varphi(\varphi,1,\varphi)$
is a direct check.  The
real-edge carriers among the
ten cyclic banks are found
by looking at which atom
interiors lie over the three
old open intervals; the
remaining isolated atomic
types are $0,2,5,8$.
The given closed atomic crests
avoid these carriers.  This
does not assert an edge-free
crest for every fine
subsector or later iterate.
\end{proof}
